\chapter{Integrators}\label{ch:integrators}

\section{The Boris push}
In LIONS units, $\vect p = \beta\gamma$, \cref{eq:lorentz} reads
\begin{equation}
\frac{\dd\vect p}{\dd t} = \frac{q}{m_0c}\,\vect E + \frac{q}{m_0\gamma}\,\vect p\times\vect B,\qquad
\frac{\dd\vect x}{\dd t} = \frac{c\,\vect p}{\gamma}.
\end{equation}
The Boris scheme, whose properties are analysed in \cite{qin2013}, keeps momenta at half steps and positions at full steps.
One step from $\vect p^{\,n-1/2}$ to $\vect p^{\,n+1/2}$ is a half electric kick, a magnetic rotation, and a second half
electric kick:
\begin{align}
\vect p^- &= \vect p^{\,n-1/2} + \frac{q\Delta t}{2m_0c}\,\vect E^n,\\
\vect p^+ - \vect p^- &= (\vect p^+ + \vect p^-)\times\vect t,\qquad \vect t = \frac{q\Delta t}{2m_0\gamma}\,\vect B^n,
\label{eq:boris-implicit}\\
\vect p^{\,n+1/2} &= \vect p^+ + \frac{q\Delta t}{2m_0c}\,\vect E^n,\\
\vect x^{n+1} &= \vect x^n + \frac{c\Delta t}{\gamma^{n+1/2}}\,\vect p^{\,n+1/2}.
\end{align}
\Cref{eq:boris-implicit} is the implicit midpoint rule for the rotation.
Dotting it with $\vect p^+ + \vect p^-$ gives $|\vect p^+| = |\vect p^-|$: the magnetic part conserves $|\vect p|$ exactly,
for any step and any field.
It has the explicit solution
\begin{equation}
\vect p' = \vect p^- + \vect p^-\times\vect t,\qquad
\vect p^+ = \vect p^- + \vect p'\times\frac{2\vect t}{1+t^2},
\end{equation}
a rotation about $\vect B$ by the angle $\alpha = 2\arctan|\vect t|$.
The exact rotation over $\Delta t$ is $\omega\Delta t$, with $|\vect t| = \omega\Delta t/2$.
LIONS writes the same scheme with a $3\times3$ matrix (\texttt{lions.pusher.boris\_push}).
It takes $\gamma$ in $\vect t$ from the previous step, which is the same thing when $\vect E = 0$.

\paragraph{Properties.}
The Boris map is second order and time-reversible.
Qin et al.\ showed that it preserves phase-space volume, although it is not symplectic \cite{qin2013}.
The distinction matters for the one-turn maps of \cref{ch:orbits}.

\section{Boris in a uniform field}\label{sec:boris-uniform}
Take $\vect E = 0$ and a uniform $B$, with exact orbit radius $\rho$ and frequency $\omega$.
Each step rotates $\vect p$ by $\alpha = 2\arctan(\omega\Delta t/2)$ and moves the particle by a chord of exact length
$v\Delta t = \rho\,\omega\Delta t$.
Consecutive chords make the angle $\alpha$, so the full-step positions lie on a circle of radius
\begin{equation}
R = \frac{\rho\,\omega\Delta t}{2\sin(\alpha/2)} = \rho\sqrt{1+\left(\frac{\omega\Delta t}{2}\right)^2},
\qquad \frac{R}{\rho} - 1 = \frac{(\omega\Delta t)^2}{8} + O(\Delta t^4),
\label{eq:boris-radius}
\end{equation}
using $\sin(\arctan u) = u/\sqrt{1+u^2}$.
The rotation per step is short by $\omega\Delta t - 2\arctan(\omega\Delta t/2) = (\omega\Delta t)^3/12 + O(\Delta t^5)$, so
the phase lags by $O(\Delta t^2)$ per turn.
\Cref{fig:boris-polygon} shows the polygon at eight steps per turn, where $R/\rho = \BorisPolyRadius$.

A drift-kick-drift variant (OPAL \texttt{LF2}) records positions half a drift after each kick.
These are the midpoints of the Boris chords, at distance $R\cos(\alpha/2) = \rho$ from the centre.
The radius test is blind for it, and case A had to add a phase metric (\cref{sec:caseA}).

\begin{figure}[t]
\centering
\includegraphics[width=0.45\textwidth]{boris_polygon}
\caption{Boris positions in a uniform field at eight steps per turn: a polygon inscribed in a circle of radius
$\rho\sqrt{1+(\omega\Delta t/2)^2}$, slightly larger than the exact orbit.}
\label{fig:boris-polygon}
\end{figure}

\section{Runge--Kutta}
The classical fourth-order Runge--Kutta method (RK4) applied to $\dot y = \lambda y$ gives $y_{n+1} = R(\lambda\Delta t)\,y_n$ with
$R(z) = 1+z+z^2/2+z^3/6+z^4/24$.
Circular motion is $\lambda = \mathrm{i}\omega$.
With $x = \omega\Delta t$,
\begin{equation}
|R(\mathrm{i}x)|^2 = \left(1-\frac{x^2}{2}+\frac{x^4}{24}\right)^2 + \left(x-\frac{x^3}{6}\right)^2
= 1 - \frac{x^6}{72} + \frac{x^8}{576},
\end{equation}
so each step shrinks the amplitude by $1-x^6/144$.
Over one turn of $n$ steps ($x = 2\pi/n$), $|\vect p|$ falls by
\begin{equation}
\frac{\Delta|\vect p|}{|\vect p|} = \frac{(2\pi)^6}{144\,n^5}.
\label{eq:rk4-decay}
\end{equation}
This is fifth order in $1/n$, while the phase error is fourth order.
The decay is secular: it adds up turn after turn (case D2, \cref{sec:caseD}).

Adaptive embedded pairs such as RKF45 (RF-Track \texttt{rkf45}) choose the step from a local error estimate.
Their error follows the requested tolerance, not the output sampling, and it also accumulates.

\section{An explicit kick is not a rotation}
A leapfrog that applies the magnetic force explicitly, $\vect p^{\,n+1} = \vect p^{\,n} + 2\,\vect p^{\,n}\times\vect t$,
adds a vector perpendicular to $\vect p$ of length $\omega\Delta t\,|\vect p|$.
So $|\vect p|^2$ grows by the factor $1+(\omega\Delta t)^2$ at every step, and $|\vect p|$ (hence the radius) by the
factor $(1+(2\pi/n)^2)^{n/2}$ per turn: a relative growth of about $2\pi^2/n$, first order in the step.
RF-Track 2.7.4 \texttt{leapfrog} showed exactly this growth on a pure magnetic field in the install checks:
$\Delta r/r = \num{5.636e-2}$ per turn at 360 steps, against the formula's \LeapfrogPred{} (\texttt{DESIGN.md},
section~7, single instrument).
It was excluded from cyclotron tracking.

\section{Symplectic and volume-preserving maps}
Let $X = (q_1,\dots,q_d,P_1,\dots,P_d)$ be canonical coordinates, and $J = \partial\Phi/\partial X$ the Jacobian of a
map $\Phi$.
The map is symplectic if
\begin{equation}
J^\mathsf{T} S J = S,\qquad S = \begin{pmatrix}0 & I_d\\ -I_d & 0\end{pmatrix}.
\label{eq:symplectic}
\end{equation}
This implies $\det J = 1$.
The converse holds only for $d = 1$, where symplectic and area-preserving are the same.
The exact flow of a Hamiltonian system is symplectic.
For a smooth Hamiltonian, a symplectic integrator exactly conserves a modified Hamiltonian close to the true one, up to
exponentially small terms, so its energy error stays bounded over very long times \cite{hairer2006}.
The argument needs smoothness; it does not cover a field with gradient jumps, such as a bilinear map.
A merely volume-preserving map has no such guarantee.

Two consequences used later:
\begin{itemize}
\item The Poincar\'e map of a Hamiltonian flow on a section at fixed energy is symplectic in the reduced coordinates.
For planar motion this is a two-dimensional map $(r,P_r)\mapsto(r',P_r')$, so $\det M_r = 1$ exactly for the exact
flow.
This holds for a continuous $B_z$ whose gradient jumps at cell edges: the variational equation $\dot J = A(t)J$ has a
trace-free $A$, and with transversal crossings of the edges the exact one-turn map stays $C^1$.
The slope changes discussed in \cref{ch:orbits} belong to the discrete map, not to the exact flow.
\item A numerical one-turn map built from a non-symplectic step, here Boris (volume-preserving in the full phase space,
but not symplectic), has no such property on the section.
Its $\det M_r - 1$ measures how far the discrete map is from the exact one.
\end{itemize}

\paragraph{Canonical momenta.}
At $z=0$, choose the gauge $A_r = 0$, $A_\theta(r,\theta) = r^{-1}\int_0^r r'B_z(r',\theta)\,\dd r'$.
Then $B_z = r^{-1}(\partial_r(rA_\theta) - \partial_\theta A_r)$, and the canonical radial momentum equals the kinetic
one, $P_r = p_r$.
Off the plane a vector potential exists only for the Maxwell-consistent field ($D = \nabla_\perp^2B$; the truncated field
with $D=0$ is divergent).
With $A_z = 0$, $A_r = (z^2/2r)\,\partial_\theta B$ and $A_\theta = r^{-1}\int_0^r r'B\,\dd r' - (z^2/2)\,\partial_rB$, the kinetic
$p_z$ is canonical and $p_r$ differs from $P_r$ by $O(z^2)$.
The change of variables has a Jacobian equal to the identity plus $O(z)$, so at small amplitude the Jacobians of
\cref{ch:orbits} can use kinetic momenta.
